\documentclass[12pt]{article}
\usepackage[T1]{fontenc}
\usepackage[utf8]{inputenc}
\usepackage{lmodern}
\usepackage{textcomp}
\usepackage{enumitem}
\usepackage{geometry}
\geometry{margin=1in}
\begin{document}
\begin{center}
Answer Key - Health Sciences - Graduate Level Assessment 18 \
\small (Answers are ordered exactly as in the question paper.)
\end{center}

\section*{Multiple Choice}
\begin{enumerate}[label=\arabic*.]
\item \textbf{Answer:} A
\\ \textit{Options:} A) The plasma concentration of vitamin B$_{12}$; B) The plasma concentration of prothrombin; C) The plasma concentration of preprothrombin; D) The plasma activity of acid phosphatase; E) The plasma activity of lipase
\\ \textit{Note:} In vitamin D deficiency, the body's ability to absorb calcium is impaired, which can disrupt the metabolism of other vitamins, including vitamin B$_{12}$, potentially leading to an increase in its plasma concentration as a compensatory mechanism.
\item \textbf{Answer:} D
\\ \textit{Options:} A) The medial wall; B) The ethmoidal wall; C) The orbital rim; D) The floor; E) The roof
\\ \textit{Note:} The floor of the orbital wall is most likely to collapse in a 'blow out' fracture due to its thinner structure and anatomical positioning, which makes it more susceptible to the increased pressure and force exerted by trauma to the eye.
\item \textbf{Answer:} A
\\ \textit{Options:} A) Coronary arteries; B) Inferior vena cava; C) Pulmonary capillaries; D) Right atrium; E) Left atrium
\\ \textit{Note:} The correct answer is the coronary arteries because they receive blood directly from the right ventricle through the pulmonary arteries, which transport deoxygenated blood to the lungs for oxygenation before returning to the heart.
\item \textbf{Answer:} A
\\ \textit{Options:} A) Determined by the individual's lifestyle habits; B) Dependent on the individual's stress levels; C) Randomly determined at birth; D) Dependent on the individual's gender; E) Dependent on the individual's diet and exercise
\\ \textit{Note:} The initial length of telomeres is influenced by an individual's lifestyle habits, such as diet, exercise, and stress levels, which can affect cellular aging and telomere maintenance.
\item \textbf{Answer:} D
\\ \textit{Options:} A) Breast Cancer; B) Influenza; C) Stroke; D) Suicide; E) Asthma
\\ \textit{Note:} Research consistently shows that men have higher rates of completed suicide compared to women, attributed to factors such as societal norms, methods of suicide chosen, and mental health stigma.
\end{enumerate}

\section*{True / False}
\begin{enumerate}[label=\arabic*.]
\setcounter{enumi}{5}
\item \textbf{Answer:} False
\\ \textit{Options:} A) True; B) False
\\ \textit{Note:} Engaging in weight lifting alone is insufficient to enhance reaction time in older adults compared to younger adults because reaction time is influenced by a combination of factors including neuromuscular coordination, flexibility, and aerobic fitness, which are not adequately addressed by weight lifting alone.
\item \textbf{Answer:} False
\\ \textit{Options:} A) True; B) False
\\ \textit{Note:} The karyotype 47,XXY, which is associated with Klinefelter syndrome, typically results in more significant developmental and health challenges compared to 47,XXX, which is associated with Triple X syndrome and generally has milder effects.
\item \textbf{Answer:} True
\\ \textit{Options:} A) True; B) False
\\ \textit{Note:} Exclusive breastfeeding for six months provides essential nutrients, antibodies, and hydration that protect newborns from waterborne illnesses and malnutrition, particularly in areas with limited access to safe water.
\item \textbf{Answer:} True
\\ \textit{Options:} A) True; B) False
\\ \textit{Note:} DNA polymerase does not require zinc ions for its catalytic activity, which is characteristic of zinc-dependent enzymes, thus confirming that it is not classified as such.
\item \textbf{Answer:} True
\\ \textit{Options:} A) True; B) False
\\ \textit{Note:} The elderly contribute to effective community change by leveraging their extensive life experience, wisdom, and established social networks to influence and mobilize others toward health initiatives and advocacy efforts.
\end{enumerate}

\section*{Long Form Response}
\begin{enumerate}[label=\arabic*.]
\setcounter{enumi}{10}
\item \textbf{Answer:} Vitamin E deficiency plays a significant role in dental health, particularly by contributing to enamel defects and increasing the risk of dental caries. Vitamin E is a powerful antioxidant that protects cell membranes from oxidative stress, which is crucial for maintaining the integrity of dental tissues. A deficiency can lead to impaired enamel formation and mineralization, resulting in weakened tooth structure and an increased susceptibility to caries. Furthermore, the lack of Vitamin E may disrupt the balance of oral microbiota, promoting pathogenic bacteria that contribute to tooth decay. Understanding these mechanisms highlights the importance of adequate Vitamin E intake in dental health and suggests that supplementation could be a beneficial strategy for preventing enamel defects and reducing caries risk.
\\ \textit{Note:} correct because it identifies Vitamin E's role as an antioxidant in protecting dental tissues, explains how its deficiency can lead to enamel defects and increased caries risk due to impaired enamel formation and mineralization, and highlights the implications for dental health.
\item \textbf{Answer:} Siblings share approximately 50\% of their genetic material due to their inheritance from the same parents, which significantly influences their human leukocyte antigen (HLA) haplotypes. HLA genes are crucial for immune system function and are inherited in blocks, meaning that siblings are likely to inherit similar combinations of these genes. Sisters, in particular, are often more likely to share a common HLA haplotype because they can inherit the same maternal and paternal alleles without the additional genetic variation introduced by male offspring. This genetic similarity not only enhances compatibility for organ transplantation but also contributes to shared immune responses. Consequently, the shared maternal lineage and the mechanisms of inheritance increase the likelihood that sisters will exhibit a common HLA haplotype compared to siblings of different sexes.
\\ \textit{Note:} correct because siblings inherit approximately 50\% of their genetic material from each parent, leading to a higher probability of sharing similar HLA haplotypes, especially among sisters who can inherit identical alleles from both the mother and father.
\end{enumerate}

\vfill
\noindent\textit{End of Answer Key}
\end{document}